\documentclass[10pt,a4paper]{amsart}
\usepackage{amsmath,amssymb,amsthm,graphicx}
\usepackage{verbatim}
\newenvironment{reference}{\comment}{\endcomment}
%\newenvironment{reference}{}{}
\newenvironment{slogan}{\comment}{\endcomment}
\newenvironment{history}{\comment}{\endcomment}
% For commutative diagrams we use Xy-pic
\usepackage[all]{xy}

% We use 2cell for 2-commutative diagrams.
\xyoption{2cell}
% Theorem environments.
%
\theoremstyle{plain}
\newtheorem{theorem}[subsection]{Theorem}
\newtheorem{proposition}[subsection]{Proposition}
\newtheorem{lemma}[subsection]{Lemma}

\theoremstyle{definition}
\newtheorem{definition}[subsection]{Definition}
\newtheorem{example}[subsection]{Example}
\newtheorem{exercise}[subsection]{Exercise}
\newtheorem{situation}[subsection]{Situation}

\theoremstyle{remark}
\newtheorem{remark}[subsection]{Remark}
\newtheorem{remarks}[subsection]{Remarks}

\numberwithin{equation}{subsection}

% Macros
%
\def\lim{\mathop{\mathrm{lim}}\nolimits}
\def\colim{\mathop{\mathrm{colim}}\nolimits}
\def\Spec{\mathop{\mathrm{Spec}}}
\def\Hom{\mathop{\mathrm{Hom}}\nolimits}
\def\Ext{\mathop{\mathrm{Ext}}\nolimits}
\def\SheafHom{\mathop{\mathcal{H}\!\mathit{om}}\nolimits}
\def\SheafExt{\mathop{\mathcal{E}\!\mathit{xt}}\nolimits}
\def\Sch{\mathit{Sch}}
\def\Mor{\mathop{\mathrm{Mor}}\nolimits}
\def\Ob{\mathop{\mathrm{Ob}}\nolimits}
\def\Sh{\mathop{\mathit{Sh}}\nolimits}
\def\NL{\mathop{N\!L}\nolimits}
\def\CH{\mathop{\mathrm{CH}}\nolimits}
\def\proetale{{pro\text{-}\acute{e}tale}}
\def\etale{{\acute{e}tale}}
\def\QCoh{\mathit{QCoh}}
\def\Ker{\mathop{\mathrm{Ker}}}
\def\Im{\mathop{\mathrm{Im}}}
\def\Coker{\mathop{\mathrm{Coker}}}
\def\Coim{\mathop{\mathrm{Coim}}}

% Boxtimes
%
\DeclareMathSymbol{\boxtimes}{\mathbin}{AMSa}{"02}

%
% Macros for moduli stacks/spaces
%
\def\QCohstack{\mathcal{QC}\!\mathit{oh}}
\def\Cohstack{\mathcal{C}\!\mathit{oh}}
\def\Spacesstack{\mathcal{S}\!\mathit{paces}}
\def\Quotfunctor{\mathrm{Quot}}
\def\Hilbfunctor{\mathrm{Hilb}}
\def\Curvesstack{\mathcal{C}\!\mathit{urves}}
\def\Polarizedstack{\mathcal{P}\!\mathit{olarized}}
\def\Complexesstack{\mathcal{C}\!\mathit{omplexes}}
% \Pic is the operator that assigns to X its picard group, usage \Pic(X)
% \Picardstack_{X/B} denotes the Picard stack of X over B
% \Picardfunctor_{X/B} denotes the Picard functor of X over B
\def\Pic{\mathop{\mathrm{Pic}}\nolimits}
\def\Picardstack{\mathcal{P}\!\mathit{ic}}
\def\Picardfunctor{\mathrm{Pic}}
\def\Deformationcategory{\mathcal{D}\!\mathit{ef}}
\usepackage{cleveref}
\crefname{originallemma}{original lemma}{original lemmas}
\Crefname{originallemma}{Original lemma}{Original lemmas}
\makeatletter
\newlabel{lemma-bijective-map}{{lemma-bijective-map}{source}{}{sourcepointer-1}{}}
\newlabel{lemma-bijective-map@cref}{{[originallemma][0][]lemma-bijective-map}{[1][0][]source}{}{}{}}
\makeatother
\crefname{originallemma}{original lemma}{original lemmas}
\Crefname{originallemma}{Original lemma}{Original lemmas}
\makeatletter
\newlabel{lemma-intersection-closed-in-quasi-compact}{{lemma-intersection-closed-in-quasi-compact}{source}{}{sourcepointer-2}{}}
\newlabel{lemma-intersection-closed-in-quasi-compact@cref}{{[originallemma][0][]lemma-intersection-closed-in-quasi-compact}{[1][0][]source}{}{}{}}
\makeatother
\crefname{originallemma}{original lemma}{original lemmas}
\Crefname{originallemma}{Original lemma}{Original lemmas}
\makeatletter
\newlabel{lemma-profinite}{{lemma-profinite}{source}{}{sourcepointer-3}{}}
\newlabel{lemma-profinite@cref}{{[originallemma][0][]lemma-profinite}{[1][0][]source}{}{}{}}
\makeatother
\crefname{originallemma}{original lemma}{original lemmas}
\Crefname{originallemma}{Original lemma}{Original lemmas}
\makeatletter
\newlabel{lemma-topological-group-limits}{{lemma-topological-group-limits}{source}{}{sourcepointer-4}{}}
\newlabel{lemma-topological-group-limits@cref}{{[originallemma][0][]lemma-topological-group-limits}{[1][0][]source}{}{}{}}
\makeatother
\title{Profinite topological groups are exactly the cofiltered limits of finite discrete topological groups}
\author{The Stacks Project Authors}
\date{stacks-project:topology.tex:local-snapshot-20261012; GNU FDL 1.2 or later}
\begin{document}
\maketitle
\noindent\textit{Editorial note.} All delivered mathematical text is the author's own, copied from the pinned Stacks Project chapter source sources/stacks-topology/pinned-original-source.tex; no statement or proof step was rewritten, reordered or removed.\par\medskip
\noindent\textit{Editorial note.} Preamble provenance: the Stacks Project keeps one single shared preamble for every chapter; this card takes it directly from the staged sources/stacks-topology/preamble.tex as three verbatim spans (lines 8-12, 14-18 and 154-217), so no derived file was written and nothing was added or redefined. The wrapper class and the base packages are supplied by the assembler.\par\medskip
\noindent\textit{Editorial note.} Bibliography provenance: neither the counted statement nor its proof contains a citation (this lane rescans the whole span before writing the card), so the card carries no bibliography and no entry was reconstructed.\par\medskip
\noindent\textit{Editorial note.} Selection provenance: this card was opened under the file-based control contract (CONTROL.md), with the mechanical pre-card gate run on its statement and proof interval before the card was written.\par\medskip
\noindent\textit{Allowed standard prerequisites (editorial).}\begin{itemize}
\item The corresponding characterization for topological spaces. Invoked at the start of the counted proof (source lines 6007-6009). Provides the limit description of the underlying space.
\item The statement that limits of topological groups are computed on the underlying topological spaces and groups. Invoked inside the counted proof (source lines 6009-6011). Carries the limit statements between spaces and groups.
\item The statement that an intersection of closed subsets of a quasi-compact space is nonempty when all finite intersections are. Invoked inside the counted proof (source line 6056). Makes the comparison map surjective.
\item The statement that a bijective map of quasi-compact Hausdorff spaces is a homeomorphism. Invoked at the end of the counted proof (source line 6057). Upgrades the isomorphism to a homeomorphism of the limit.
\end{itemize}\medskip
\ifdefined\hypertarget\hypertarget{sourcepointer-1}{}\fi
\noindent\textit{Source reference.} Original lemma lemma-bijective-map: sources/stacks-topology/pinned-original-source.tex, source lines 3498--3498. Navigation pointer only: the counted text refers to the Stacks item lemma-bijective-map, which is an existing labeled result of this same chapter and is not part of this card. The pointer carries the target's own label as its number and its definition line as its span; its content is not reproduced here.\par\medskip
\ifdefined\hypertarget\hypertarget{sourcepointer-2}{}\fi
\noindent\textit{Source reference.} Original lemma lemma-intersection-closed-in-quasi-compact: sources/stacks-topology/pinned-original-source.tex, source lines 1941--1941. Navigation pointer only: the counted text refers to the Stacks item lemma-intersection-closed-in-quasi-compact, which is an existing labeled result of this same chapter and is not part of this card. The pointer carries the target's own label as its number and its definition line as its span; its content is not reproduced here.\par\medskip
\ifdefined\hypertarget\hypertarget{sourcepointer-3}{}\fi
\noindent\textit{Source reference.} Original lemma lemma-profinite: sources/stacks-topology/pinned-original-source.tex, source lines 4351--4351. Navigation pointer only: the counted text refers to the Stacks item lemma-profinite, which is an existing labeled result of this same chapter and is not part of this card. The pointer carries the target's own label as its number and its definition line as its span; its content is not reproduced here.\par\medskip
\ifdefined\hypertarget\hypertarget{sourcepointer-4}{}\fi
\noindent\textit{Source reference.} Original lemma lemma-topological-group-limits: sources/stacks-topology/pinned-original-source.tex, source lines 5975--5975. Navigation pointer only: the counted text refers to the Stacks item lemma-topological-group-limits, which is an existing labeled result of this same chapter and is not part of this card. The pointer carries the target's own label as its number and its definition line as its span; its content is not reproduced here.\par\medskip
\begin{lemma}
\label{lemma-profinite-group}
Let $G$ be a topological group. The following are equivalent
\begin{enumerate}
\item $G$ as a topological space is profinite,
\item $G$ is a limit of a diagram of finite discrete topological groups,
\item $G$ is a cofiltered limit of finite discrete topological groups.
\end{enumerate}
\end{lemma}

\begin{proof}
We have the corresponding result for topological spaces, see
Lemma \ref{lemma-profinite}. Combined with
Lemma \ref{lemma-topological-group-limits}
we see that it suffices to prove that (1) implies (3).

\medskip\noindent
We first prove that every neighbourhood $E$ of the neutral element $e$ contains
an open subgroup. Namely, since $G$ is the cofiltered limit of finite
discrete topological spaces (Lemma \ref{lemma-profinite}),
we can choose a continuous map $f : G \to T$ to a finite discrete
space $T$ such that $f^{-1}(f(\{e\})) \subset E$.
Consider
$$
H = \{g \in G \mid f(gg') = f(g')\text{ for all }g' \in G\}
$$
This is a subgroup of $G$ and contained in $E$. Thus it suffices to
show that $H$ is open. Pick $t \in T$ and set $W = f^{-1}(\{t\})$.
Observe that $W \subset G$ is open and closed, in particular quasi-compact.
For each $w \in W$ there exist open neighbourhoods $e \in U_w \subset G$
and $w \in U'_w \subset W$ such that $U_wU'_w \subset W$.
By quasi-compactness we can find $w_1, \ldots, w_n$ such that
$W = \bigcup U'_{w_i}$. Then
$U_t = U_{w_1} \cap \ldots \cap U_{w_n}$ is an open neighbourhood
of $e$ such that $f(gw) = t$ for all $w \in W$.
Since $T$ is finite we see that $\bigcap_{t \in T} U_t \subset H$
is an open neighbourhood of $e$. Since $H \subset G$ is a subgroup
it follows that $H$ is open.

\medskip\noindent
Suppose that $H \subset G$ is an open subgroup. Since $G$ is quasi-compact
we see that the index of $H$ in $G$ is finite. Say
$G = Hg_1 \cup \ldots \cup Hg_n$. Then
$N = \bigcap_{i = 1, \ldots, n} g_iHg_i^{-1}$ is
an open normal subgroup contained in $H$. Since $N$ also has finite
index we see that $G \to G/N$ is a surjection to a finite discrete
topological group.

\medskip\noindent
Consider the map
$$
G \longrightarrow \lim_{N \subset G\text{ open and normal}} G/N
$$
We claim that this map is an isomorphism of topological groups.
This finishes the proof of the lemma as the limit on the right
is cofiltered (the intersection of two open normal subgroups is
open and normal). The map is continuous as each $G \to G/N$ is continuous.
The map is injective as $G$ is Hausdorff and every neighbourhood
of $e$ contains an $N$ by the arguments above.
The map is surjective by
Lemma \ref{lemma-intersection-closed-in-quasi-compact}.
By Lemma \ref{lemma-bijective-map} the map is a homeomorphism.
\end{proof}

\end{document}
